\documentclass[11pt]{article}
\usepackage[a5paper,margin=16mm]{geometry}
\usepackage{fontspec}
\setmainfont{TeX Gyre Pagella}
\setsansfont{TeX Gyre Heros}
\usepackage{graphicx}
\usepackage{xcolor}
\usepackage{caption}
\usepackage{subcaption}
\usepackage{float}
\usepackage{rotating}
\usepackage{mwe}
\captionsetup{font=small,labelfont=bf,labelsep=period,justification=raggedright,singlelinecheck=false}
\captionsetup[sub]{font=footnotesize}

\begin{document}
\listoffigures
\bigskip

\begin{figure}[H]
  \centering
  \begin{subfigure}[t]{.45\linewidth}
    \centering
    \includegraphics[width=\linewidth]{example-image-a}
    \caption{First panel}\label{fig:a}
  \end{subfigure}\hfill
  \begin{subfigure}[t]{.45\linewidth}
    \centering
    \includegraphics[width=\linewidth,angle=90,origin=c,height=.6\linewidth,keepaspectratio]{example-image-b}
    \caption{Rotated second panel}\label{fig:b}
  \end{subfigure}
  \caption{Two example images placed side by side; panels~\subref{fig:a} and~\subref{fig:b} share a baseline.}\label{fig:pair}
\end{figure}

Figure~\ref{fig:pair} contains \ref{fig:a} and \ref{fig:b}.
\begin{figure}[H]
  \centering
  \scalebox{0.5}[0.35]{\fbox{\includegraphics[width=.5\linewidth]{example-image-c}}}
  \caption{A scaled, framed image}\label{fig:scaled}
\end{figure}
\begin{table}[H]
  \centering
  \caption{Inline table beside nothing}\label{tab:t}
  \begin{tabular}{lr}\hline Item & Qty\\\hline nuts & 12\\ bolts & 40\\\hline\end{tabular}
\end{table}

\begin{sidewaysfigure}
  \centering
  \includegraphics[width=.6\linewidth]{example-image-16x9}
  \caption{Sideways figure with a 16:9 image}\label{fig:wide}
\end{sidewaysfigure}
See Figures~\ref{fig:scaled} and~\ref{fig:wide} and Table~\ref{tab:t}.
\end{document}
